\documentclass[conference]{IEEEtran}
\IEEEoverridecommandlockouts
% The preceding line is only needed to identify funding in the first footnote. If that is unneeded, please comment it out.
\usepackage{cite}
\usepackage{amsmath,amssymb,amsfonts}
\usepackage{algorithmic}
\usepackage{graphicx}
\usepackage{textcomp}
\usepackage{xcolor}
\usepackage{booktabs}
\usepackage{array}
\usepackage{multirow}
\usepackage{url}
\usepackage{eurosym}
\usepackage{textcomp}
\usepackage{hyperref}

\def\BibTeX{{\rm B\kern-.05em{\sc i\kern-.025em b}\kern-.08em
    T\kern-.1667em\lower.7ex\hbox{E}\kern-.125emX}}
    
% ===================== ADDED THIS =====================

\renewcommand{\thesection}{\arabic{section}}
\renewcommand{\thesubsection}{\thesection.\arabic{subsection}}
\renewcommand{\thesubsectiondis}{\thesubsection}

% =====================================================

\begin{document}

\title{Deep Learning for Panoramic Dental Radiographs: A Comprehensive
Survey on Caries Detection, Periodontal Bone Loss Assessment and
Explainable AI%

}

\author{
\IEEEauthorblockN{C S Deepak}
\IEEEauthorblockA{\textit{Dept.\ of Computer Science}\\
\textit{PES University -- RR Campus}\\
Bangalore, India}
\and
\IEEEauthorblockN{Chennupati Gunadeep}
\IEEEauthorblockA{\textit{Dept.\ of Computer Science}\\
\textit{PES University -- RR Campus}\\
Bangalore, India}
\and
\IEEEauthorblockN{Archana}
\IEEEauthorblockA{\textit{Dept.\ of Computer Science}\\
\textit{PES University -- RR Campus}\\
Bangalore, India}
\and
\IEEEauthorblockN{AVSN Sai Srujan}
\IEEEauthorblockA{\textit{Dept.\ of Computer Science}\\
\textit{PES University -- RR Campus}\\
Bangalore, India}
\and
\IEEEauthorblockN{\centering Under the guidance of Dr.\ Mamatha H R, Director, Bajaj Engineering Skills Training @PESU, Bangalore, India}
}
\maketitle

% -----------------------------------------------------------------------
\begin{abstract}
Dental panoramic radiographs (OPGs) are widely used for diagnosing
dental caries and periodontal disease, yet manual interpretation remains
time-consuming and prone to inter-examiner variability. Despite rapid
advances in deep learning, existing research remains fragmented, with
most studies addressing isolated tasks and lacking clinically integrated
solutions.

This paper presents a comprehensive survey of deep learning methods for
panoramic dental radiograph analysis, covering caries detection, severity
staging, periodontal bone loss assessment, tooth segmentation, and
explainable AI. A total of 47 studies published between 2021 and 2026
are systematically analyzed across architectures including CNNs,
YOLO-based detectors, U-Net variants, and transformer-based models. This
paper formally defines a unified multi-task dental AI framework
integrating detection, localization, severity staging, quantification,
and explainability within a single system.

Our analysis shows that while individual tasks achieve high performance,
no existing approach provides a complete, clinically usable pipeline.
Five key gaps are identified, including task fragmentation, absence of
severity staging, limited explainability adoption, dataset bias, and
weak clinical integration.

This work establishes a structured foundation for developing
next-generation, clinically deployable dental AI systems.
\end{abstract}

\begin{IEEEkeywords}
Deep Learning, Dental Caries, Periodontal Disease, Panoramic
Radiographs, Tooth Segmentation and Detection, Multi-Task Learning,
Explainable AI
\end{IEEEkeywords}

% -----------------------------------------------------------------------
\section{Introduction}
\label{sec:intro}

Dental panoramic radiography, also known as OrthoPantomography (OPG),
is the standard imaging tool in routine dental practice. A single
panoramic exposure captures the full dentition, surrounding alveolar
bone, and anatomical landmarks in one image, making it the primary
modality for detecting dental caries, periodontal bone loss, impacted
teeth, and other oral pathologies. Despite this clinical centrality,
accurate interpretation of panoramic radiographs demands substantial
expertise and is inherently subject to inter-examiner variability,
particularly among early-stage practitioners such as dental students
and interns, who represent the population most vulnerable to diagnostic
error.

Dental caries, defined as the progressive bacterial demineralization of
tooth enamel and dentine, affects approximately 2.3 billion people
globally and is the most prevalent non-communicable disease worldwide
\cite{luke2025accuracy}. Periodontal disease, characterized by
irreversible alveolar bone loss around the teeth, affects over 19\% of
the adult global population and is the leading cause of adult tooth
loss \cite{khubrani2025detection}. Both conditions share a critical
clinical property: their early radiographic signs are subtle, easily
missed, and disproportionately consequential when detection is delayed.
A faint radiolucency at the proximal surface of a premolar, or a
1--2\,mm reduction in alveolar bone height, can be clinically decisive
yet visually ambiguous even for trained examiners.

Over the past five years, deep learning has enabled significant advances
in automated dental radiograph analysis. Convolutional neural networks
(CNNs), YOLO-based object detectors, U-Net segmentation architectures,
and transformer-based models have demonstrated strong performance across
individual diagnostic tasks such as caries detection, tooth
segmentation, and periodontal bone loss estimation
\cite{nassiri2025yolo,jundaeng2025advanced,bhoopalan2026self}.

However, despite these advances, current research remains fundamentally
fragmented at the system level. Existing approaches predominantly focus
on isolated tasks or single-disease objectives, without integrating
these components into a unified diagnostic workflow required in clinical
practice. This reflects a deeper structural limitation: current models
are optimized for task-specific performance rather than holistic
diagnostic reasoning \cite{zanini2024systematic,muresanu2024automating}.

The challenge of automated dental radiograph analysis is inherently
multi-dimensional. A clinically viable system must simultaneously
address multi-instance detection across all teeth, identification of
low-contrast early-stage lesions, precise spatial measurement of
anatomical landmarks such as the cemento-enamel junction (CEJ) and
alveolar bone crest, and the coexistence of multiple diseases within a
single panoramic image. This multi-stage diagnostic process can be
conceptualized as a pipeline, as illustrated in Fig.~\ref{fig:pipeline}.

\begin{figure*}[tp]
    \centering
    \includegraphics[width=\textwidth]{fig1.jpeg}
    \caption{System Pipeline of Proposed AI Model.}
    \label{fig:pipeline}
\end{figure*}

Furthermore, the lack of standardized datasets and evaluation protocols
limits reproducibility and cross-study comparability. A significant
proportion of studies rely on private, single-center datasets, leading
to dataset bias and poor generalization to external clinical
environments \cite{muresanu2024automating,delamare2024panoramic}. In
addition, explainable artificial intelligence (XAI) remains severely
underutilized, despite its critical role in clinical trust,
interpretability, and regulatory acceptance
\cite{ferrara2025applications,bilal2024pdcnet}. The absence of
explainability, combined with inconsistent validation practices, creates
a substantial barrier to real-world deployment.

Unlike prior surveys that focus on individual tasks or model
architectures, this work adopts a system-level perspective, analyzing
how disparate components can be integrated into a unified diagnostic
framework \cite{zhou2025deep,albano2024artificial}.

This survey reviews 47 research papers published between 2021 and 2026,
providing a comprehensive synthesis of deep learning methods applied to
dental radiograph analysis. The literature is organized into six
functional categories: caries detection and severity classification,
periodontal disease and bone loss analysis, multi-task end-to-end
diagnosis, explainable AI, dataset benchmarking and clinical validation,
and tooth segmentation and numbering. The central finding of this survey
is that no existing system simultaneously addresses caries detection,
severity staging, periodontal bone loss quantification, and
explainability within a single unified framework for panoramic
radiographs.

The main contributions of this survey are as follows:
\begin{itemize}
    \item A task-level taxonomy of deep learning systems for panoramic
          dental radiograph analysis, covering detection, segmentation,
          classification, and multi-task architectures.
    \item A formal problem formulation for unified multi-disease dental
          diagnosis, highlighting the integration requirements of
          detection, severity staging, quantification, and
          explainability.
    \item A cross-category comparative analysis of 47 studies,
          identifying inconsistencies in datasets, evaluation metrics,
          and validation protocols.
    \item A quantified gap analysis revealing key limitations in current
          research, including single-disease model focus, absence of
          automated severity staging, dataset fragmentation, and low
          adoption of explainable AI.
    \item A research roadmap toward clinically deployable, explainable,
          multi-task dental AI systems.
\end{itemize}

This survey establishes unified, explainable, multi-task dental AI as
the critical next step toward clinically deployable diagnostic systems,
particularly for supporting junior clinicians and dental students in
real-world practice.




% -----------------------------------------------------------------------
\section{Search Methodology}
\label{sec:methodology}

This survey used a structured keyword search across five digital
libraries: IEEE Xplore, PubMed/MEDLINE, Scopus, Google Scholar, and
SpringerLink. The search was executed between January 2024 and March
2025, with publication year filters restricting results to 2020 onward.

\subsection{Search Strategy}

Search queries were organized into five thematic clusters:

\begin{itemize}
  \item \textit{Caries detection on panoramic radiographs:}
        ``dental caries'' panoramic radiograph deep learning; ``caries
        detection'' panoramic radiograph CNN; ``caries segmentation''
        orthopantomogram neural network.

  \item \textit{Periodontal disease and bone loss:} ``periodontal
        disease'' panoramic radiograph deep learning; ``alveolar bone
        loss'' orthopantomogram AI; ``CEJ'' ``bone crest'' panoramic
        radiograph segmentation.

  \item \textit{Multi-disease and end-to-end OPG systems:}
        ``multi-task'' dental panoramic radiograph deep learning;
        ``dental disease detection'' panoramic radiograph end-to-end.

  \item \textit{Explainable AI in dental imaging:} ``explainable AI''
        dental radiograph; ``Grad-CAM'' dental panoramic;
        ``interpretable deep learning'' dental imaging.

  \item \textit{Dataset and clinical validation:} ``dental panoramic
        radiograph dataset'' deep learning; ``OPG benchmark'' AI;
        ``clinical validation'' dental AI panoramic.
\end{itemize}

\subsection{Inclusion and Exclusion Criteria}

Papers were included if they: (i)~applied deep learning or machine
learning to dental radiograph analysis; (ii)~used panoramic, bitewing,
or periapical radiographs as primary input; (iii)~targeted detection,
segmentation, classification, or staging of dental caries, periodontal
disease, or related pathologies; and (iv)~reported quantitative
performance metrics. Papers were excluded if they focused exclusively on
CBCT/3D imaging, addressed only non-dental imaging, or did not report
reproducible quantitative results. Clinical reviews without AI models
were included only when they provided foundational clinical definitions
directly relevant to the survey scope.

\subsection{Final Corpus}

After applying the criteria, 47~papers were selected, spanning 2021--2026,
with 25~papers from 2025 (53\%), 17~from 2024 (36\%), 2~from 2026,
2~from 2023, and 1~from 2021 -- reflecting the field's rapid recent
growth. Publication venues include BMC Oral Health (8~papers), IEEE (8),
MDPI journals (9), Elsevier (4), Springer (3), and Scientific
Reports/Nature Portfolio (3).

% -----------------------------------------------------------------------
\section{System-Level Taxonomy and Problem Formulation}
\label{sec:taxonomy}

\subsection{Formal Problem Definition}

Panoramic dental radiograph analysis is inherently a multi-component
diagnostic problem that cannot be reduced to single-task prediction.
In clinical practice, meaningful diagnosis requires coordinated
execution of multiple interdependent processes: anatomical localization
of teeth, identification of pathological conditions, assessment of
disease severity, and quantitative measurement of clinically relevant
parameters. Most existing deep learning systems in the reviewed
literature address only one component in isolation
\cite{albano2024artificial,luke2025accuracy,zanini2024systematic},
producing outputs that are diagnostically incomplete and clinically
difficult to integrate.

A unified dental AI system maps an input OPG to structured diagnostic
outputs through six core functional components: detection, localization,
severity staging, quantification, explainability, and clinical
integration. To the best of our knowledge, no reviewed study in the
corpus of 47 papers fully addresses all six components simultaneously
on panoramic radiographs, representing the primary architectural gap
motivating this survey.

These six core components are defined as follows:

\begin{itemize}
    \item \textbf{Detection:} Identification of pathological conditions
          from radiographic features (carious lesions, periodontal bone
          loss, periapical pathology, impacted teeth). YOLO-family
          architectures dominate this task
          \cite{ayhan2025detection,nassiri2025yolo,chen2024cariesxrays,
          bizjak2024advancing,pornprasertsuk2025clinical}.

    \item \textbf{Localization:} Spatial identification of affected
          teeth through bounding boxes or polygon segmentation.
          Tooth-level localization using FDI numbering is present in
          pipeline systems
          \cite{jundaeng2025advanced,chen2024cariesxrays,bizjak2024advancing,
          pornprasertsuk2025clinical} but absent in
          classification-only approaches
          \cite{kaushik2025densenet201,mathur2025comparative,ozdemir2026enhanced}.

    \item \textbf{Severity Staging:} Classification of disease
          progression along clinically meaningful scales. For caries:
          enamel, dentine, pulp involvement
          \cite{asadi2024detection,bizjak2024advancing}. For
          periodontal disease: Stages I--IV per 2018 AAP/EFP
          classification \cite{chmielewski20252018}. This is the most
          critically underaddressed component, with only two C1 papers
          providing caries severity differentiation
          \cite{bizjak2024advancing,asadi2024detection}.

    \item \textbf{Quantification:} Measurement of the CEJ--alveolar
          bone crest distance and percentage bone loss. Only three
          studies address this: Jundaeng et al.\ achieved
          mAP50\,=\,0.995 \cite{jundaeng2025advanced}; Wimalasiri et
          al.\ reported ICC\,=\,0.80 \cite{wimalasiri2025ai}; and Tian
          et al.\ achieved ICC\,=\,0.83 \cite{tian2025development}.

    \item \textbf{Explainability:} Generation of interpretable visual
          outputs for clinical verification. Only 3 of 47 papers
          implement explainability (6.4\%): Chen et al.\
          \cite{chen2023deep} and Bilal et al.\ \cite{bilal2024pdcnet}
          use Grad-CAM, and Chen et al.\ \cite{chen2024cariesxrays}
          show feature map attention.

    \item \textbf{Clinical Integration:} Alignment with diagnostic
          workflows, reporting standards, and regulatory requirements.
          Only Schwendicke et al.\ \cite{schwendicke2021cost} model
          economic impact, and only Lee et al.\ \cite{lee2025evaluation}
          measure clinician acceptance through a 65-dentist usability
          survey.
\end{itemize}

This six-component formulation stands in direct contrast to existing
approaches that address each task in isolation -- a fragmentation
independently identified as critical by multiple systematic reviews
\cite{iacob2025automated,albano2024artificial,bonfanti2025deep,
khubrani2025detection,zanini2024systematic,zhou2025deep}.

\subsection{System Taxonomy}

Existing approaches can be categorized into four distinct system
archetypes based on task integration and diagnostic scope. This taxonomy
reveals both progression toward integration and persistent limitations
preventing full clinical utility. Table~\ref{tab:taxonomy} summarizes
capability coverage across taxonomy dimensions.

\subsubsection{Single-Task Approaches}

Single-task systems address one diagnostic objective in isolation --
caries detection, tooth segmentation, or periodontal classification --
without integration. This accounts for the majority of reviewed studies
(18 C1, 15 C2 papers). While several achieve high performance --
CariesXrays FPCL reaching AP50\,=\,84.1\% \cite{chen2024cariesxrays},
Attention U-Net achieving Recall\,=\,0.96
\cite{pornprasertsuk2025clinical}, U-Net variants reporting
Dice\,=\,96.77\% \cite{vickram2025extraction} -- none produce severity
grading, periodontal assessment, or clinical explanations. The
fundamental limitation is not performance but scope: optimizing one
component without modeling interdependence with other diagnostic tasks
produces clinically incomplete systems. Systematic reviews confirm this
\cite{albano2024artificial,zanini2024systematic}, identifying absence
of multi-task integration as the primary barrier to clinical
translation.

\subsubsection{Multi-Task Approaches}

Seven reviewed papers adopt multi-task designs with varying
completeness. Chen et al.\ \cite{chen2023deep} is the most clinically
complete, classifying both periodontitis and caries on a per-tooth
basis using EfficientNet-B0 with YOLOv7 localization (ROC
AUC\,=\,98.67\% for periodontitis, 98.31\% for caries) with Grad-CAM
explainability. However, this system operates exclusively on periapical
radiographs -- not panoramic OPGs -- limiting direct applicability.
Multi-task approaches expose a critical performance asymmetry:
structural tasks (segmentation, enumeration) consistently achieve
near-ceiling performance, while pathological tasks (disease detection,
severity staging) remain substantially harder. Mure\c{s}anu et al.\
\cite{muresanu2024automating} quantified the real-world cost: YOLOv8
F1 dropped from 0.60 (internal) to 0.47 (external) -- a 22\%
degradation from single-center training.

\subsubsection{Pipeline-Based Approaches}

Pipeline systems decompose diagnosis into sequential stages, reflecting
the staged clinical reasoning dentists employ. Jundaeng et al.\
\cite{jundaeng2025advanced} implemented a precise periodontal pipeline:
YOLOv8 Stage~1 for tooth segmentation (mAP50\,=\,0.92) feeds Stage~2
for CEJ and bone crest detection (mAP50\,=\,0.995,
Sensitivity\,=\,1.0). Pornprasertsuk-Damrongsri et al.\
\cite{pornprasertsuk2025clinical} combined YOLOv5s for tooth detection
with Attention U-Net for caries segmentation (Recall\,=\,0.96,
F1\,=\,0.85). The inherent limitation is error propagation:
inaccuracies in early stages compound through subsequent stages. None
of the reviewed pipelines address severity staging, and explainability
appears in only one system \cite{chen2023deep} -- operating on
periapical, not panoramic, radiographs.

\subsubsection{End-to-End Approaches}

End-to-end models directly map radiographs to diagnostic outputs using
deep architectures -- primarily transformers. Bhoopalan et al.\
\cite{bhoopalan2026self} applied Vision Transformer and Swin Transformer
to 5,098 OPGs, achieving ViT accuracy\,=\,96\% without spatial
localization or severity grading. Van Nistelrooij et al.\
\cite{vannistelrooij2024combining} combined Mask DINO with Swin-B
Transformer, achieving tooth-level assessment performance
(FROC\,=\,0.988) but with 40-second inference time unsuitable for
clinical use. The fundamental barrier is data: the field lacks
large-scale, multi-disease, multi-center public datasets required for
end-to-end learning.

The system-level taxonomy and capability gaps across different architectural paradigms are illustrated in Fig.~\ref{fig:taxonomy}

\begin{figure*}[t]
  \centering
  \includegraphics[width=0.85\textwidth]{fig4.jpeg}
  \caption{System-level taxonomy of dental AI approaches for panoramic radiographs, 
  illustrating four architectural paradigms (Single-Task, Multi-Task, Pipeline-Based, 
  and End-to-End Systems) and their respective capability coverage across detection, 
  localization, severity, quantification, and explainability (XAI) dimensions, 
  alongside the proposed Unified System as a future research direction.}
  \label{fig:taxonomy}
\end{figure*}

\begin{table*}[t]
\renewcommand{\arraystretch}{1.5}
\caption{System Taxonomy --- Representative Studies by Design Pattern
and Capability Coverage}
\label{tab:taxonomy}
\centering
\scriptsize
\setlength{\tabcolsep}{4pt}
\begin{tabular}{p{1.6cm} p{3.5cm} p{1.0cm} p{1.2cm} p{1.2cm} p{1.0cm} p{1.0cm} p{1.5cm}}
\toprule
\textbf{System Type} & \textbf{Representative Studies} &
\textbf{Detect} & \textbf{Localize} & \textbf{Severity} &
\textbf{Quant.} & \textbf{XAI} & \textbf{Multi-Disease} \\
\midrule
Single-Task &
\cite{ayhan2025detection,alharbi2023detection,khuwaja2024ai,
chen2024cariesxrays,perez2024ai,zancan2025ai,ozdemir2026enhanced}
& \checkmark & Partial & \texttimes & \texttimes & \texttimes & \texttimes \\

Multi-Task &
\cite{nassiri2025yolo,chen2023deep,gamal2025automated,
vannistelrooij2024combining,muresanu2024automating,bhoopalan2026self}
& \checkmark & \checkmark & \texttimes & \texttimes & Partial & Partial \\

Pipeline-Based &
\cite{jundaeng2025advanced,chen2024cariesxrays,bizjak2024advancing,
pornprasertsuk2025clinical}
& \checkmark & \checkmark & \texttimes & Partial & \texttimes & \texttimes \\

End-to-End &
\cite{vannistelrooij2024combining,bhoopalan2026self}
& \checkmark & Partial & \texttimes & \texttimes & \texttimes & Partial \\

\midrule
\textbf{Unified (Required)} & \textit{None in reviewed corpus}
& \checkmark & \checkmark & \checkmark & \checkmark & \checkmark & \checkmark \\
\bottomrule
\multicolumn{8}{l}{\footnotesize \checkmark\,=\,fully addressed;\
\texttimes\,=\,not addressed;\ Partial\,=\,partially addressed.}\\
\multicolumn{8}{l}{\footnotesize Quant.\,=\,quantitative measurement;\
XAI\,=\,explainable AI outputs.}\\
\end{tabular}
\end{table*}

\subsection{Task Hierarchy}

Dental radiograph analysis follows a structured sequence of
interdependent processing stages reflecting the actual cognitive process
trained clinicians employ:

\begin{figure*}[t]
  \centering
  \includegraphics[width=0.95\textwidth]{fig5.jpeg}
  \caption{Task hierarchy for panoramic dental radiograph analysis, illustrating 
  the sequential stages from anatomical structure extraction to clinical report 
  generation. Early stages such as tooth segmentation and numbering achieve high 
  performance, while downstream stages including severity assessment, quantification, 
  and explainability remain underexplored, highlighting key limitations in current 
  systems.}
  \label{fig:taskhierarchy}
\end{figure*}

Early stages focus on anatomical structure identification, where systems
perform well: tooth segmentation achieves pooled sensitivity\,=\,0.92,
specificity\,=\,0.94 \cite{bonfanti2025deep}; tooth enumeration reaches
mAP50\,=\,0.978 \cite{nassiri2025yolo}. Middle stages -- disease
detection and severity assessment -- represent the performance
bottleneck. Disease detection on OPG reaches only mAP50\,$\approx$\,0.61
\cite{nassiri2025yolo}, and severity staging is attempted in only two
of 18 C1 papers \cite{bizjak2024advancing,asadi2024detection}. Terminal
stages -- explainability and structured reporting -- are most neglected:
explainability appears in only 3 of 47 papers (6.4\%)
\cite{chen2023deep,chen2024cariesxrays,bilal2024pdcnet}, and structured
report generation is absent from all systems. This hierarchy
demonstrates that dental AI is fundamentally a sequential reasoning
problem, where failures in early stages propagate to downstream
diagnostic decisions. For systems targeting junior clinicians, the
absence of visual explanations eliminates educational value entirely.

\subsection{System-Level Limitations}

Five system-level limitations persist across the reviewed literature,
reflecting field-wide methodological choices:

\begin{itemize}
    \item \textbf{Task Fragmentation:} No panoramic OPG system
          simultaneously addresses caries detection, severity staging,
          periodontal quantification, and explainability. Chen et al.\
          \cite{chen2023deep} achieves dual-disease detection with
          explainability -- closest to unified -- but only on periapical
          radiographs.

    \item \textbf{Severity Staging Deficit:} Only Asadi
          \cite{asadi2024detection} and Bizjak and Robic
          \cite{bizjak2024advancing} provide caries severity
          differentiation among 18 C1 papers. For periodontal disease,
          the 2018 AAP/EFP Stages I--IV \cite{chmielewski20252018} are
          not addressed by any system.

    \item \textbf{Explainability Absence:} Only 3 of 47 papers
          implement explainability (6.4\%), contradicting field
          consensus
          \cite{nassiri2025yolo,gamal2025automated,sastika2025ai,
          ferrara2025applications,khubrani2025detection}. The EU AI Act
          and FDA guidance require transparency for high-risk systems --
          requirements 43 of 47 systems cannot satisfy.

    \item \textbf{Dataset Fragmentation:} Approximately 60\% of studies
          use private, single-center datasets causing heterogeneity
          values $I^2$\,=\,88--96\%
          \cite{iacob2025automated,albano2024artificial,bonfanti2025deep}.
          Mure\c{s}anu et al.\ \cite{muresanu2024automating} documented
          22\% F1 degradation from internal to external data. No
          large-scale, multi-disease, multi-center, publicly available
          panoramic dataset with standardized annotations exists.

    \item \textbf{Clinical Workflow Disconnection:} Except for
          Schwendicke et al.\ \cite{schwendicke2021cost}
          (cost-effectiveness analysis) and Lee et al.\
          \cite{lee2025evaluation} (65-dentist usability study), no
          system evaluates AI within actual clinical workflows. Most
          produce bounding boxes without structured reports or treatment
          recommendations.
\end{itemize}
%---------------------------

\section{Background}
\label{sec:background}

\subsection{Panoramic Radiography in Clinical Dentistry}

Orthopantomography (OPG) produces a two-dimensional tomographic
representation of the curved dental arches on a flat image plane. Unlike
periapical or bitewing radiographs that show specific regions, a single
panoramic image captures all 32 permanent teeth, maxilla and mandible,
and surrounding structures -- making it the most information-dense
single radiographic acquisition in routine dental practice. For caries
detection, panoramic radiographs are less sensitive than bitewing for
early proximal lesions \cite{perez2024ai,bizjak2024advancing,pornprasertsuk2025clinical}.
For periodontal bone loss assessment, OPG is the standard screening
tool, enabling visualization of the alveolar crest relative to the
cemento-enamel junction (CEJ) across all teeth
\cite{jundaeng2025advanced,iacob2025automated,khubrani2025detection}.
This modality distinction directly shapes the architectural requirements
for AI systems addressing each task.

\subsection{Deep Learning Architectures}

The architectures reviewed fall into four functional families. First,
CNNs with transfer learning -- ResNet~\cite{khuwaja2024ai},
DenseNet~\cite{kaushik2025densenet201,kaushik2025densenet201,ozdemir2026enhanced},
VGG~\cite{mathur2025comparative}, EfficientNet~\cite{chen2023deep},
Inception~\cite{gamal2025automated,zancan2025ai} -- are used primarily
for classification, leveraging ImageNet pre-training to compensate for
limited dental dataset sizes. Second, YOLO-family one-stage object
detectors (YOLOv5--YOLOv9)
\cite{ayhan2025detection,nassiri2025yolo,bizjak2024advancing,muresanu2024automating,pornprasertsuk2025clinical}
dominate detection tasks for their real-time speed and strong bounding
box localization. Third, U-Net and nested variants
\cite{alharbi2023detection,asci2024deep,lim2024dual,kong2024integrating,zanini2024systematic}
are the standard for pixel-level segmentation, using skip connections to
preserve spatial detail. Fourth, transformer architectures -- ViT
\cite{bhoopalan2026self}, Swin Transformer
\cite{vannistelrooij2024combining,zhang2025reliable}, RT-DETR
\cite{karamuftuo2025panoramic} -- represent the most recent frontier,
offering global attention at higher computational cost.

\subsection{Evaluation Metrics}

A persistent challenge across the reviewed literature is metric
inconsistency that prevents direct cross-study comparison
\cite{iacob2025automated,bonfanti2025deep,khubrani2025detection,zanini2024systematic}.
Object detection studies report mAP at IoU thresholds of 0.50 and
0.50:0.95. Segmentation studies report IoU and Dice coefficient.
Classification studies report accuracy, AUC, sensitivity, specificity,
and F1. Several studies \cite{chen2023deep,bilal2024pdcnet,zancan2025ai}
additionally report PPV and NPV -- metrics with direct clinical
relevance that standard accuracy obscures. The absence of a standardized
dental AI evaluation protocol is identified as one of the field's most
critical methodological gaps
\cite{iacob2025automated,luke2025accuracy,khubrani2025detection}.

% -----------------------------------------------------------------------
\section{Caries Detection and Severity Classification (C1)}
\label{sec:caries}

Caries detection is the most extensively studied task in dental AI,
accounting for 18 of the 47 reviewed papers. Studies span object
detection, segmentation, classification, and systematic review
methodologies across panoramic, bitewing, periapical, and intraoral RGB
modalities.

\subsection{YOLO-Based Detection on Panoramic Radiographs}

Ayhan et al.~\cite{ayhan2025detection} proposed a two-stage pipeline
targeting caries beneath fixed dental prostheses (FDPs) -- a clinically
challenging scenario where metal artifacts from crowns and bridges
actively obscure lesion visibility. YOLOv7 detected FDP regions
(mAP\,=\,0.968, F1\,=\,0.956) in Stage~1; a CBAM-enhanced YOLOv7
applied to cropped ROIs achieved caries detection with mAP\,=\,0.846
and F1\,=\,0.830 in Stage~2. The CBAM attention module demonstrated
measurable benefit for low-contrast lesion boundaries beneath metallic
artifacts.

Bizjak and Robic~\cite{bizjak2024advancing} applied a two-step YOLOv8
pipeline on the public DENTEX dataset, achieving tooth enumeration
mAP@0.5\,=\,0.95 and lesion sensitivity\,=\,0.70 (internal), 0.67
(external) -- one of the few studies providing genuine external clinical
validation, which revealed the performance degradation characteristic of
single-center training. Asadi~\cite{asadi2024detection} applied Faster
R-CNN to the UFBA-UESC public dataset, achieving 94\% overall accuracy
but reporting a critical clinical weakness: recall for mild
caries\,=\,0.64 versus severe caries\,=\,0.69 -- the opposite of
clinical need, where early lesions are most important to detect.
Pornprasertsuk-Damrongsri et al.~\cite{pornprasertsuk2025clinical}
combined YOLOv5s for tooth detection with Attention U-Net for caries
segmentation on 500 panoramic radiographs with bitewing-confirmed ground
truth, achieving Recall\,=\,0.96, F1\,=\,0.85 with Bland-Altman
agreement analysis.An example of multi-tooth carious lesion detection with severity differentiation in panoramic radiographs is shown in Fig.~\ref{fig:dentex} .

\begin{figure*}[t]
    \centering
    \includegraphics[width=0.65\textwidth]{fig2.jpg}
    \caption{Example panoramic radiograph with annotated carious lesions, 
    illustrating multi-instance detection and severity differentiation 
    (caries vs.\ deep caries) across multiple teeth.}
    \label{fig:dentex}
\end{figure*}

\subsection{U-Net Segmentation Approaches}

Alharbi et al.~\cite{alharbi2023detection} compared U-Net, U-Net++, and
U-Net3+ on 1,500 panoramic images. U-Net3+ achieved the best test
performance (IoU\,=\,0.67, Dice\,=\,0.60, Accuracy\,=\,0.95), with its
full-scale skip connections providing measurable advantage over standard
U-Net (Dice\,=\,0.58). The gap between training IoU (0.77) and test IoU
(0.67) reveals overfitting from severe class imbalance -- only 510 of
1,500 images contained visible cavities.

Lim et al.~\cite{lim2024dual} proposed Faster R-CNN + U-Net
(EfficientNet-B0 encoder) on 597 panoramic radiographs, achieving
IoU\,=\,0.608, Dice\,=\,0.743, Recall\,=\,0.731 on 33,752 augmented
cropped tooth images, though without severity grading or explainability.
Asci et al.~\cite{asci2024deep} applied U-Net to 6,075 pediatric
panoramic radiographs stratified across primary, mixed, and permanent
dentition, achieving F1\,=\,0.8675, Sensitivity\,=\,0.8269,
Precision\,=\,0.9123 -- demonstrating U-Net's adaptability across
diverse tooth morphologies, though AUC\,=\,0.76 remains below the 90\%
clinical threshold. Vickram and Mathan Muthu~\cite{vickram2025extraction}
applied UNet variants on 500 dental X-ray images, achieving
Accuracy\,=\,98.03\% and Dice\,=\,96.77\%, though on a notably small
dataset with no clinical validation.


\subsection{CNN Classification and Transfer Learning}

Zancan et al.~\cite{zancan2025ai} applied InceptionResNet-v2 to 7,304
tooth-level images from 935 panoramic radiographs, achieving
accuracy\,=\,0.94 and F1\,=\,0.94 with five-fold cross-validation --
one of the stronger validation approaches in C1. Kaushik and
Sharma~\cite{kaushik2025densenet201} applied DenseNet201 to 29,815
dental X-ray images, achieving 91\% overall accuracy but a
cavity-specific F1 of only 0.27 -- a stark illustration of class
imbalance consequences. Ozdemir et al.~\cite{ozdemir2026enhanced}
applied DenseNet-121 to 14,498 cropped tooth images (2,303 caries /
12,195 non-caries), achieving accuracy\,=\,93.17\% and F1\,=\,87.49\%,
though without localization on full panoramic images. Khuwaja et
al.~\cite{khuwaja2024ai} evaluated YOLOv8, ResNet-50, EfficientDet, and
MobileNetV2 on 7,447 intraoral RGB smartphone images -- targeting mobile
diagnostics for resource-limited settings -- with ResNet-50 achieving
mAP@0.50\,=\,0.670.

\subsection{Systematic Reviews and Meta-Analyses}

Albano et al.~\cite{albano2024artificial} reviewed 20 studies via
PRISMA/QUADAS-2, finding accuracy 73\%--98\% and AUC 0.84--0.98 across
CNN-based models, with high heterogeneity and an explicit call for
standardized datasets and multicenter validation. Luke and
Rezallah~\cite{luke2025accuracy} conducted a meta-analysis of 21 studies
(17,190 total samples), reporting the best individual study achieving
98.8\% accuracy but $I^2$ values of 88--96\% indicating severe
heterogeneity that undermines pooled estimates. Zanini et
al.~\cite{zanini2024systematic} reviewed 42 studies confirming that most
works are single-task and lack caries severity differentiation as a
systematic field-wide gap. Sastika and Marappan~\cite{sastika2025ai}
reviewed accuracy ranges of 85--99\% across multiple architectures,
notably acknowledging Grad-CAM as a requirement in surveyed studies.

Chen et al.~\cite{chen2024cariesxrays} introduced the CariesXrays
dataset (6,000 hospital-grade panoramic X-rays, 13,783 expert-verified
bounding boxes) alongside the FPCL architecture -- achieving
AP\,=\,48.2\%, AP50\,=\,84.1\%, outperforming YOLOv8
(AP50\,=\,80.7\%). CariesXrays is the most significant public benchmark
contribution in the C1 category.

\subsection{Bitewing Caries and Cost-Effectiveness}

Perez de Frutos et al.~\cite{perez2024ai} evaluated RetinaNet,
YOLOv5-M, and EfficientDet on 13,887 bitewing X-rays from the HUNT4
Oral Health Study, annotated by six dentists with a 197-image consensus
test set. YOLOv5-M achieved mAP\,=\,0.647, F1\,=\,0.548 -- matching or
surpassing individual expert annotators, providing strong evidence that
deep learning performs at human expert level on proximal caries
detection. Schwendicke et al.~\cite{schwendicke2021cost} combined U-Net
performance evaluation (accuracy\,=\,0.80, sensitivity\,=\,0.75 vs.\
dentist sensitivity\,=\,0.36) with a Markov lifetime cost-effectiveness
simulation on 3,686 bitewing radiographs, projecting increased tooth
retention (64 vs.\ 62 years) and reduced cost per tooth
(\euro298 vs.\ \euro322), yielding ICER\,=\,$-$13.9~\euro/year --
establishing AI caries detection as a cost-dominant strategy.

\begin{table*}[t]
\caption{Summary of Reviewed Studies --- Caries Detection and Severity Classification (C1)}
\label{tab:c1}
\centering
\small
\begin{tabular}{p{0.6cm} p{4.8cm} p{3.0cm} p{3.5cm} p{1.0cm} p{3.2cm}}
\toprule
\textbf{Ref} & \textbf{Model} & \textbf{Dataset Size} & \textbf{Best Metric} & \textbf{XAI} & \textbf{Key Gap} \\
\midrule
\cite{ayhan2025detection}   & YOLOv7 + CBAM & 1,004 OPG & mAP=0.846, F1=0.830 & None & Single-center, no XAI \\
\cite{khuwaja2024ai}        & YOLOv8, ResNet-50, EfficientDet, MobileNetV2 & 7,447 RGB & mAP@0.50=0.670 & None & RGB only, no radiographs \\
\cite{alharbi2023detection} & U-Net, U-Net++, U-Net3+ & 1,500 OPG & IoU=0.67, Dice=0.60 & None & Class imbalance, no severity \\
\cite{asci2024deep}         & U-Net & 6,075 OPG (pediatric) & F1=0.8675, AUC=0.76 & None & Below 90\% clinical threshold \\
\cite{albano2024artificial} & Systematic Review (CNNs) & 20 studies & Acc up to 0.98, AUC 0.98 & None & High heterogeneity \\
\cite{kaushik2025densenet201} & DenseNet201 & 29,815 X-rays & Acc=91\%, Cavity F1=0.27 & None & Severe class imbalance \\
\cite{chen2024cariesxrays}  & FPCL (Faster R-CNN + D2D-FPN + Contrastive) & 6,000 OPG & AP=48.2\%, AP50=84.1\% & Attn maps & No severity/periodontal \\
\cite{perez2024ai}          & YOLOv5-M, RetinaNet, EfficientDet & 13,887 bitewing & mAP=0.647, F1=0.548 & None & Bitewing only, no periodontal \\
\cite{sastika2025ai}        & Review: U-Net, ResNet, YOLOv7/v8, ensemble & Multiple studies & Acc 85--99\%, F1 87--94\% & Discussed & No original experiment \\
\cite{schwendicke2021cost}  & U-Net & 3,686 bitewing & Acc=0.80, ICER=$-$\euro13.9/yr & None & Single-center, German context \\
\cite{luke2025accuracy}     & Meta-analysis (CNNs) & 21 studies, 17,190 images & Best Acc=98.8\% & None & $I^2$=88--96\% heterogeneity \\
\cite{vickram2025extraction}& UNet, UNet+, UNet++ & 500 images & Acc=98.03\%, Dice=96.77\% & None & Small dataset, no clinical val. \\
\cite{zancan2025ai}         & InceptionResNet-v2 & 935 OPG / 7,304 teeth & Acc=0.94, F1=0.94 & None & No localization, no periodontal \\
\cite{lim2024dual}          & Faster R-CNN + U-Net (EfficientNet-B0) & 597 OPG & IoU=0.608, Dice=0.743 & None & No severity, no periodontal \\
\cite{zanini2024systematic} & Systematic Review (U-Net, Mask R-CNN, ResNet) & 42 studies & Various (review) & None & No severity differentiation \\
\cite{bizjak2024advancing}  & YOLOv8 two-step pipeline & 1,005 OPG + 85 external & mAP@0.5=0.95, Sens=0.70 & None & Caries only, no severity \\
\cite{asadi2024detection}   & Faster R-CNN & 1,000 OPG & Acc=94\%, Mild recall=0.64 & None & Weak mild caries detection \\
\cite{pornprasertsuk2025clinical} & YOLOv5s + Attention U-Net & 500 OPG, 14,997 teeth & Recall=0.96, F1=0.85 & None & Single-disease, no periodontal \\
\cite{ozdemir2026enhanced}  & DenseNet-121 & 14,498 tooth images & Acc=93.17\%, F1=87.49\% & None & No localization, no severity \\
\bottomrule
\end{tabular}
\end{table*}

% -----------------------------------------------------------------------
\section{Periodontal Disease and Alveolar Bone Loss Detection (C2)}
\label{sec:perio}

Periodontal disease detection from radiographs requires quantifying the
spatial relationship between anatomical landmarks -- specifically the
distance between the cemento-enamel junction (CEJ) and the alveolar
bone crest -- across all teeth in the panoramic image. This demands
precise spatial localization and continuous measurement rather than
binary classification. Fifteen papers in the C2 category reflect a less
mature but rapidly advancing field.

\subsection{Clinical Foundation}

Chmielewski et al.~\cite{chmielewski20252018} provide essential clinical
framing: the 2018 AAP/EFP classification defines periodontitis
stages~I--IV and grades~A--C based on bone loss percentage, probing
depth, tooth loss, and risk factors. Across 11 clinical studies covering
459 cases and 852 examiners, staging accuracy peaked at 90\% but
inter-examiner agreement remained variable -- the subjectivity problem
that AI systems can address, but only if they produce quantitative,
staged outputs aligned with the 2018 criteria rather than binary bone
loss labels. Fig.~\ref{fig:perio} illustrates the spatial nature of periodontal bone loss assessment, requiring accurate identification of anatomical landmarks such as the cemento-enamel junction and alveolar bone crest.

\begin{figure*}[t]
    \centering
    \includegraphics[width=0.65\textwidth]{fig3.jpg}
    \caption{Example panoramic radiograph illustrating periodontal bone 
    loss regions. Accurate assessment requires spatial measurement between 
    the cemento-enamel junction and alveolar bone crest.}
    \label{fig:perio}
\end{figure*}

\subsection{Landmark Detection and Bone Loss Quantification}

Jundaeng et al.~\cite{jundaeng2025advanced} presented the most
technically complete periodontal AI system reviewed. Their dual YOLOv8
pipeline on 2,000 panoramic radiographs achieved tooth segmentation
(Accuracy\,=\,0.97, mAP50\,=\,0.92) and CEJ/bone level detection
(Accuracy\,=\,0.98, Sensitivity\,=\,1.0, Specificity\,=\,0.98,
F1\,=\,0.90, mAP50\,=\,0.995). Sensitivity\,=\,1.0 for bone level
detection is clinically important -- missing true bone loss is more
harmful than a false positive in a screening context. Limitations
include single-center data, high-quality images only, and no caries
integration or explainability.

Iacob et al.~\cite{iacob2025automated} synthesized evidence from 20
diagnostic accuracy studies covering 141,166 radiographs through
PRISMA/QUADAS-2, finding that U-Net, YOLO, ResNet, ConvNeXt, and
DeNTNet have all been applied but high heterogeneity prevents reliable
pooling. Khubrani et al.~\cite{khubrani2025detection} meta-analyzed 30
studies using APPRAISE-AI, reporting pooled sensitivity\,=\,87\%,
specificity\,=\,76\%, accuracy\,=\,84\% -- useful but insufficient for
autonomous clinical use.

Mathur et al.~\cite{mathur2025comparative} applied VGG-16 to 1,874
panoramic radiographs for binary bone loss classification
(accuracy\,=\,89.27\%, sensitivity\,=\,83.27\%) -- the most common C2
approach but also the most clinically limited, providing no staging
information. Cerda Mardini et al.~\cite{cerda2024determining} attempted
quantitative staging on 500 panoramic radiographs, reporting
F1\,=\,0.29 for light-moderate bone loss and F1\,=\,0 for severe bone
loss -- honestly documenting the difficulty of quantitative staging and
the model's failure on the most clinically critical class. Kong et
al.~\cite{kong2024integrating} integrated SRGAN super-resolution
preprocessing with U-Net for bone loss segmentation
(Dice\,=\,0.91, IoU\,=\,84.9\%) -- the strongest segmentation metrics
in C2 -- demonstrating that preprocessing quality directly impacts
segmentation accuracy, though the dataset of 698 images is small and no
disease staging is provided.

Wimalasiri et al.~\cite{wimalasiri2025ai} applied YOLOv8x and Keypoint
R-CNN to 1,000 intraoral periapical X-rays, achieving ICC\,=\,0.80 for
severity and 87\% for bone loss pattern classification (horizontal
vs.\ vertical) -- the only study addressing pattern classification,
which has direct prognostic implications for surgical planning. Tian et
al.~\cite{tian2025development} developed the PAN-POL mathematical model
validated on 290 panoramic patients, achieving ICC\,=\,0.83, though
this is a rule-based approach rather than a deep learning system.

\subsection{Periodontitis Classification}

Bilal et al.~\cite{bilal2024pdcnet} proposed PDCNET -- a custom
five-block CNN with Grad-CAM explainability -- compared against
EfficientNet-B0, DenseNet-201, VGG-16, VGG-19, Inception-V3, and
MobileNet on a Kaggle dataset augmented with SMOTE-Tomek, achieving
accuracy\,=\,98.39\%, AUC\,=\,99.79\%, F1\,=\,98.31\%. These are the
highest classification metrics in C2, though the 50+50 original sample
size raises serious concerns about statistical meaningfulness. The paper
is valuable for demonstrating Grad-CAM feasibility in periodontal AI.
Ferrara et al.~\cite{ferrara2025applications} reviewed 13 periodontal
deep learning studies via PRISMA/QUADAS-2, reporting accuracy ranges of
73\%--98.6\% and noting the complete absence of explainable AI across
reviewed systems. Lee et al.~\cite{lee2025evaluation} evaluated AI
alveolar bone loss measurement usability through a survey of 65 dentists
on 550 bitewing images -- one of the only clinical usability studies in
the entire reviewed corpus.

Kim et al.~\cite{kim2024unsupervised} proposed an unsupervised few-shot
learning architecture combining UNet, CVAE, and clustering methods
(k-means/DBSCAN/GMM) on 1,372 panoramic radiographs, achieving a 14\%
improvement over supervised ViT, demonstrating the potential of
label-efficient learning in data-scarce dental AI settings.
Karamuftuo\u{g}lu et al.~\cite{karamuftuo2025panoramic} benchmarked
YOLOv12x, RT-DETR-L, and RT-DETR-X on a 387-image pediatric panoramic
dataset for furcation lesion detection, achieving mAP@0.5\,=\,0.434 --
reflecting both the difficulty of the task and the limitation of small
single-center datasets.

\begin{table*}[t]
\caption{Summary of Reviewed Studies --- Periodontal Disease and Bone Loss (C2)}
\label{tab:c2}
\centering
\small
\begin{tabular}{p{0.6cm} p{4.5cm} p{3.0cm} p{3.8cm} p{1.0cm} p{2.8cm}}
\toprule
\textbf{Ref} & \textbf{Model} & \textbf{Dataset Size} & \textbf{Best Metric} & \textbf{XAI} & \textbf{Key Gap} \\
\midrule
\cite{jundaeng2025advanced}  & Dual YOLOv8 (teeth + CEJ/bone) & 2,000 OPG & mAP50=0.995, Sens=1.0 & None & Single-center, no caries \\
\cite{chmielewski20252018}   & Clinical review (no ML) & 459 cases, 852 examiners & Stage acc=90\% & None & No AI system proposed \\
\cite{iacob2025automated}    & Systematic Review (CNN, YOLO, U-Net) & 141,166 radiographs & Various (review) & None & High heterogeneity \\
\cite{mathur2025comparative} & VGG-16 (transfer learning) & 1,874 OPG & Acc=89.27\% & None & Binary only, no staging \\
\cite{khubrani2025detection} & Meta-analysis (CNN, YOLO, U-Net, SVM) & 30 studies & Sens=87\%, Spec=76\% & None & Pooled heterogeneity \\
\cite{tian2025development}   & PAN-POL mathematical model & 290 PAN patients & ICC=0.83 & None & Not deep learning, manual \\
\cite{lee2025evaluation}     & Faster R-CNN + DeepLab v3+ & 550 bitewing & Acc=82--87\% (severe) & None & Bitewing only, no OPG \\
\cite{cerda2024determining}  & Hybrid DCNN (Xception) + rule-based PBL & 500 OPG & F1=0.29 (lt-mod), F1=0 (severe) & None & Fails on severe PBL \\
\cite{wimalasiri2025ai}      & YOLOv8x + Keypoint R-CNN & 1,000 IOPA & ICC=0.80, Pattern Acc=87\% & None & Periapical only, no OPG \\
\cite{kong2024integrating}   & SRGAN + U-Net & 698 OPG & Dice=0.91, IoU=84.9\% & None & No severity, small dataset \\
\cite{ferrara2025applications}& Systematic Review (CNNs, hybrid) & 13 studies & Acc 73--98.6\% & None & No XAI in any reviewed study \\
\cite{bilal2024pdcnet}       & PDCNET (custom CNN) + Grad-CAM & 50+50 orig (SMOTE aug.) & Acc=98.39\%, AUC=99.79\% & Grad-CAM & Dataset too small \\
\cite{karamuftuo2025panoramic} & YOLOv12x, RT-DETR-L, RT-DETR-X & 387 OPG (pediatric) & mAP@0.5=0.434 & None & Small dataset, modest mAP \\
\cite{kim2024unsupervised}   & UNet + CVAE + k-means/DBSCAN/GMM & 1,372 OPG & +14\% over supervised ViT & None & No bounding boxes \\
\bottomrule
\end{tabular}
\end{table*}

% -----------------------------------------------------------------------
\section{Multi-Task and End-to-End OPG Diagnosis (C3)}
\label{sec:multitask}

Seven papers address multiple diagnostic tasks within a single pipeline
applied to panoramic radiographs -- the architectural direction most
aligned with real clinical utility.

Chen et al.~\cite{chen2023deep} present the most clinically complete
dual-disease system reviewed, combining YOLOv7 for tooth localization
with EfficientNet-B0 for simultaneous multi-label classification of both
periodontitis and dental caries on periapical X-rays, with Grad-CAM
visualization. EfficientNet-B0 achieved ROC AUC\,=\,98.67\% for
periodontitis and 98.31\% for dental caries with 10-fold
cross-validation. The critical limitation is modality: periapical, not
panoramic.

Nassiri and Akhloufi~\cite{nassiri2025yolo} systematically compared YOLO
v5--v9 on DENTEX across quadrant segmentation, tooth enumeration, and
disease detection. YOLOv9e-seg achieved mAP50\,=\,0.978 for tooth
enumeration and 0.959 for segmentation, but pathology detection peaked
at mAP50\,=\,0.61 -- confirming that structural tasks are reliably
solvable while pathology detection remains the harder problem.

Van Nistelrooij et al.~\cite{vannistelrooij2024combining} combined
OdontoAI and DENTEX datasets with Mask DINO (ViT backbone) and Swin-B
Transformer, achieving FROC\,=\,0.988 for tooth segmentation and
AUC\,=\,0.996 for impacted teeth -- the highest-performing tooth-level
assessment reviewed. The 40-second inference time per image is the
primary barrier to clinical deployment. Gamal et
al.~\cite{gamal2025automated} demonstrated that a Weighted Average
Ensemble of CNN architectures (Custom CNN, ResNet-50, InceptionV3,
DenseNet121, VGG19) achieves accuracy\,=\,92.93\% on 1,814 cropped
tooth images, with ESRGAN super-resolution preprocessing contributing
approximately 3\% accuracy improvement.

Mure\c{s}anu et al.~\cite{muresanu2024automating} provided the most
methodologically honest external validation in C3: YOLOv8 tested on a
Romanian clinical dataset showed internal F1\,=\,0.60 degrading to
external F1\,=\,0.47 -- a 22\% performance drop that directly quantifies
the generalization gap between controlled training and real multi-center
clinical data. Bhoopalan et al.~\cite{bhoopalan2026self} applied ViT and
Swin Transformer to 5,098 OPG images, achieving ViT accuracy\,=\,96\%,
precision\,=\,95.8\%, recall\,=\,96.2\% -- strong results for
image-level classification, though without localization or segmentation.
Zhou et al.~\cite{zhou2025deep} provided a comprehensive systematic
review of datasets, methodologies, and emerging challenges in dental
image analysis, reinforcing the field's critical need for standardization.

%----------------------------------------------------------
\section{Explainable AI in Dental Diagnostics (C4)}
\label{sec:xai}

Of the 47 papers reviewed, only three report any form of explainability
output: Chen et al.~\cite{chen2023deep} use Grad-CAM for dual-disease
classification on periapical X-rays; Bilal et al.~\cite{bilal2024pdcnet}
apply Grad-CAM to periodontal disease classification; and Chen et
al.~\cite{chen2024cariesxrays} show attention visualization in feature
maps for caries detection. This constitutes \textbf{6.4\%} of the
reviewed corpus -- directly quantifying the explainability gap in dental
AI.

Gradient-weighted Class Activation Mapping (Grad-CAM) generates a
coarse localization map highlighting image regions most influential for a
CNN prediction. In the dental diagnostic context, Grad-CAM serves three
clinical functions: (i)~it allows clinicians to verify that the model
attends to anatomically plausible regions; (ii)~it provides a visual
teaching tool for junior clinicians learning radiographic interpretation;
and (iii)~it creates an audit trail relevant to the EU AI Act and FDA
guidance on AI/ML-based medical devices -- both of which classify AI
diagnostic tools as high-risk systems requiring transparency and human
oversight \cite{chen2023deep}.

Despite a clear consensus across the literature that explainability is
required for clinical adoption
\cite{nassiri2025yolo,gamal2025automated,sastika2025ai,ferrara2025applications},
the implementation rate remains at 6.4\%. This gap is particularly acute
for systems targeting junior clinicians and dental students, who require
not just a prediction but a visual justification that teaches them to
recognize the radiographic pattern and builds diagnostic confidence. The
complete absence of SHAP, LIME, or transformer attention visualization in
any reviewed dental AI study represents an untapped methodological
opportunity -- transformer attention maps are intrinsic to the
architecture and require no additional computation
\cite{bhoopalan2026self}.


%------------------------------------------------------------------

\section{Comparative Analysis of Representative Methods}
\label{sec:comparison}

\begin{table*}[t]
\renewcommand{\arraystretch}{1.4}
\caption{Comparison of Representative Deep Learning Methods for Dental
Radiograph Analysis. Performance metrics are reported as provided in
the respective studies and are not directly comparable due to variations
in tasks, datasets, and evaluation protocols.}
\label{tab:comparison}
\centering
\scriptsize
\setlength{\tabcolsep}{3pt}
\begin{tabular}{p{0.5cm} p{2.0cm} p{1.3cm} p{2.2cm} p{1.1cm} p{2.4cm} p{1.6cm} p{0.6cm} p{1.3cm} p{2.2cm}}
\toprule
\textbf{Ref} & \textbf{Task} & \textbf{Scope} &
\textbf{Dataset (Size)} & \textbf{Modality} & \textbf{Model} &
\textbf{Best Metric} & \textbf{XAI} &
\textbf{Clinical Val.} & \textbf{Key Limitation} \\
\midrule
\cite{ayhan2025detection}        & Caries Detection             & Single-task & Private OPG (1,004)             & OPG      & YOLOv7 + CBAM           & mAP=0.846        & No      & Internal & No external validation \\
\cite{jundaeng2025advanced}      & Bone Loss + Periodontal      & Multi-task  & Private OPG (2,000)             & OPG      & YOLOv8                  & mAP50=0.995      & No      & Internal & Single-center dataset \\
\cite{chen2023deep}              & Multi-disease Detection      & Multi-task  & Periapical (11k)                & X-ray    & YOLOv7 + EfficientNet   & AUC$\approx$98\% & Yes     & Internal & Not panoramic \\
\cite{alharbi2023detection}      & Caries Segmentation          & Single-task & DNS (1,500)                     & OPG      & U-Net3+                 & Dice=0.60        & No      & Internal & Small dataset \\
\cite{asci2024deep}              & Caries Segmentation          & Single-task & Pediatric OPG (6,075)           & OPG      & U-Net                   & F1=0.867         & No      & Internal & No severity staging \\
\cite{gamal2025automated}        & Multi-disease Classification & Multi-task  & Roboflow (1,075)                & OPG      & CNN Ensemble            & Acc=92.9\%       & No      & Internal & No explainability \\
\cite{chen2024cariesxrays}       & Caries Detection             & Single-task & CariesXrays (6,000)             & OPG      & Faster R-CNN            & AP50=84.1\%      & Partial & Internal & Single-disease \\
\cite{perez2024ai}               & Caries Detection             & Single-task & HUNT4 (13k)                     & Bitewing & YOLOv5                  & mAP=0.647        & No      & External & Not panoramic \\
\cite{nassiri2025yolo}           & Multi-task OPG Analysis      & Multi-task  & DENTEX ($\sim$1,000)            & OPG      & YOLOv9-seg              & mAP=0.978        & No      & Internal & Limited pathology \\
\cite{vannistelrooij2024combining}& Tooth + Diagnosis            & Multi-task  & OdontoAI (1,661) + DENTEX (682) & OPG      & Mask DINO + Swin        & mAP=0.848        & No      & Internal & Slow pipeline \\
\cite{zancan2025ai}              & Tooth + Caries               & Multi-task  & 935 PAN + 18k teeth             & OPG      & InceptionResNet         & Acc=0.998        & No      & Internal & No localization \\
\cite{lim2024dual}               & Caries Segmentation          & Single-task & 597 PAN                         & OPG      & Faster R-CNN + U-Net    & Dice=0.743       & No      & Internal & Single-disease \\
\cite{bizjak2024advancing}       & Caries Detection             & Single-task & DENTEX + external               & OPG      & YOLOv8                  & mAP=0.95         & No      & External & No segmentation \\
\cite{asadi2024detection}        & Caries Severity              & Single-task & 1,000 PAN                       & OPG      & Faster R-CNN            & Acc=94\%         & No      & Internal & Weak mild detection \\
\cite{pornprasertsuk2025clinical}& Caries Detection             & Multi-stage & 500 PAN                         & OPG      & YOLOv5 + U-Net          & Recall=0.96      & No      & Internal & Single-disease \\
\cite{mathur2025comparative}     & Bone Loss Classification     & Single-task & 1,874 OPG                       & OPG      & VGG-16                  & Acc=89.27\%      & No      & Internal & No localization \\
\cite{lee2025evaluation}         & Bone Loss Detection          & Single-task & Bitewing (550)                  & Bitewing & Faster R-CNN + DeepLab  & Acc$\approx$85\% & No      & Internal & Limited modality \\
\cite{wimalasiri2025ai}          & Bone Loss Severity           & Single-task & DenPAR (1,000)                  & IOPA     & YOLOv8 + KRCNN          & ICC=0.80         & No      & Internal & Not panoramic \\
\cite{kong2024integrating}       & Bone Loss Segmentation       & Single-task & CBNUH (100) + Kaggle ($\sim$600)& OPG      & SRGAN + U-Net           & Dice=0.91        & No      & Internal & No detection \\
\cite{muresanu2024automating}    & Multi-disease Detection      & Single-task & Multi-center panoramic          & OPG      & YOLOv8                  & F1=0.60→0.47     & No      & External & Poor generalization \\
\cite{zhang2025reliable}         & Tooth Segmentation           & Multi-stage & Tufts + O2PR                    & OPG      & Mask R-CNN + YOLO       & mAP=77.4         & No      & Internal & No disease labels \\
\cite{bonfanti2025deep}          & Tooth Segmentation Review    & Review      & 20 studies                      & Multi    & CNNs                    & Sens=0.92        & No      & Review   & High heterogeneity \\
\cite{khubrani2025detection}     & Bone Loss Review             & Review      & 30 studies                      & Multi    & CNNs                    & Acc=84\%         & No      & Review   & No standardization \\
\cite{zhou2025deep}              & Deep Learning Survey         & Review      & 260 studies                     & Multi    & Various                 & ---              & No      & Review   & No unified system \\
\cite{bilal2024pdcnet}           & Periodontal Classification   & Single-task & Kaggle ($\sim$1,000)            & X-ray    & PDCNET                  & Acc=98.39\%      & Yes     & Internal & No localization \\
\end{tabular}
\end{table*}
% -----------------------------------------------------------------------
\section{Datasets, Benchmarks, and Clinical Validation (C5)}
\label{sec:datasets}

Dataset availability and quality represent the single most frequently
cited limitation across all 47 reviewed papers. The key public
benchmarks in the field are as follows:

\begin{itemize}
  \item \textbf{DENTEX}: 1,005 panoramic radiographs with quadrant,
        tooth, and pathology annotations from three institutions
        \cite{nassiri2025yolo,vannistelrooij2024combining,bizjak2024advancing}.
  \item \textbf{CariesXrays}: 6,000 hospital-grade panoramic images
        with 13,783 caries bounding boxes \cite{chen2024cariesxrays}.
  \item \textbf{Tufts Dental Database}: Approximately 1,000 panoramic
        radiographs with tooth segmentation masks
        \cite{cerda2024determining,kim2024unsupervised,zhang2025reliable,yaswanth2025automated}.
  \item \textbf{HUNT4 Oral Health Study}: 13,887 bitewing X-rays with
        six-dentist annotations and a 197-image consensus set
        \cite{perez2024ai}.
  \item \textbf{UFBA-UESC}: 4,000 panoramic images
        \cite{asadi2024detection}.
\end{itemize}

Despite these contributions, approximately 60\% of reviewed primary
studies rely exclusively on private, single-center institutional
datasets. Systematic reviews
\cite{iacob2025automated,albano2024artificial,luke2025accuracy,bonfanti2025deep,khubrani2025detection}
consistently report $I^2$ heterogeneity values of 88--96\% as the
primary barrier to reliable evidence synthesis. The performance
degradation from internal to external validation quantified by
Mure\c{s}anu et al.~\cite{muresanu2024automating} -- F1 dropping from
0.60 to 0.47 -- directly measures the real cost of single-center
training.

Waqas et al.~\cite{waqas2025synthetic} addressed data scarcity through
GAN-based synthetic OPG generation (DCGAN, WGAN) on 4,777 real images,
producing synthetic panoramic radiographs with expert realism assessed
above 50\% relative to real images -- a promising augmentation approach,
though downstream disease detection improvement was not demonstrated.
Ameli et al.~\cite{ameli2025automated} applied YOLOv8 to 1,000 OPGs for
automated quality assessment (contrast/density accuracy\,=\,97.9\%,
artifact accuracy\,=\,87.2\%, overall quality\,=\,79.3\%) --
establishing automated quality filtering as a viable preprocessing gate
for diagnostic AI pipelines. Delamare et al.~\cite{delamare2024panoramic}
documented heterogeneous quality criteria across 41 panoramic ML
studies, calling for standardized quality protocols as a prerequisite for
reproducible AI development.

% -----------------------------------------------------------------------
\section{Tooth Segmentation and Numbering (C6)}
\label{sec:segmentation}

Four papers address tooth detection, segmentation, and numbering --
enabling upstream tasks that are prerequisites for tooth-specific disease
attribution. Accurate FDI-standard tooth numbering is required to report
that caries is in tooth~26 rather than merely somewhere in the panoramic
image.

Bonfanti-Gris et al.~\cite{bonfanti2025deep} synthesized 20 tooth
detection and segmentation studies in a PRISMA/QUADAS-2 meta-analysis,
reporting pooled sensitivity\,=\,0.92 and specificity\,=\,0.94 with
DOR\,=\,186 -- strong aggregate evidence, though $I^2 \approx 96\%$
confirms severe heterogeneity. Zhang et al.~\cite{zhang2025reliable}
proposed a multi-stage system (MobileNetV3 + SADP-Net Mask R-CNN + Swin
Transformer + YOLOv5) across three public datasets, achieving
AP50\,=\,97.3\% on Tufts and 98.1\% on O2PR for tooth instance
segmentation. Yaswanth et al.~\cite{yaswanth2025automated} applied
Attention-integrated U-Net++ to the Tufts dataset, achieving
mAP\,=\,96.4\% and Dice\,=\,90\%. Adnan and
Umer~\cite{adnan2024orthopantomogram} contributed the Orthopantomogram
Teeth Segmentation and Numbering Dataset -- 250 expert-annotated OPGs on
Zenodo -- explicitly noting the absence of disease annotations and
calling for expansion to caries and periodontal labeling.

\section{Cross-Study Trends and Observations}

This section presents a synthesized analysis of recurring patterns identified across the 47 reviewed studies. Based on the global comparison in Table~\ref{tab:comparison}, key trends emerge across model design, dataset characteristics, task complexity, and deployment challenges that collectively define the current landscape of dental AI.

\subsection{Model Architecture Trends}

\subsubsection{Transformer-Based vs. CNN-Based Architectures}

Transformer-based architectures, such as Vision Transformer (ViT) and Swin Transformer, demonstrate strong capability in capturing global contextual information, particularly for classification and multi-label diagnosis tasks~\cite{vannistelrooij2024combining}. However, as observed in Table~\ref{tab:comparison}, CNN-based architectures dominate the majority of studies, especially in detection and pipeline-based systems~\cite{kaushik2025densenet201}. This dominance is primarily due to their efficiency, lower computational requirements, and robustness on moderately sized dental datasets.

Transformer-based models typically require large annotated datasets and exhibit higher inference latency, limiting their practicality in clinical environments~\cite{vannistelrooij2024combining}. Consequently, despite their representational power, CNN-based approaches remain more suitable for real-world deployment under current dataset constraints.

\subsubsection{Dominance of YOLO in Detection Tasks}

A clear trend across the reviewed literature is the widespread adoption of YOLO-based architectures for detection-related tasks. As reflected in Table~\ref{tab:comparison}, YOLO variants are the most frequently used models for caries localization, tooth detection, and bone loss identification~\cite{nassiri2025yolo}. Their dominance stems from the ability to achieve high detection accuracy in a single inference pass while maintaining near real-time performance, making them particularly suitable for clinical applications.

\subsubsection{Ensemble Methods}

Ensemble methods improve classification robustness by combining multiple models; however, they introduce increased computational complexity and reduced inference efficiency, limiting their applicability in real-time clinical settings~\cite{gamal2025automated}.

\subsection{Data and Generalization}

Dataset size positively influences model performance, but it is not sufficient to ensure generalization. Large datasets may still suffer from issues such as class imbalance, leading to poor task-specific performance~\cite{kaushik2025densenet201}. More importantly, models trained on single-center datasets consistently exhibit reduced performance when evaluated on external data.

This generalization gap is a structural limitation across the field. For example, models trained on single-center panoramic datasets demonstrate significant performance degradation under multi-center conditions~\cite{muresanu2024automating}. Systematic reviews further confirm high heterogeneity across studies due to differences in imaging protocols, annotation standards, and model architectures~\cite{khubrani2025detection}. These findings indicate that dataset diversity, multi-center representation, and standardized evaluation protocols are essential for developing clinically reliable systems.

\subsection{Task Complexity and Learning Paradigm}

\subsubsection{Structural Tasks vs. Disease Detection}

The reviewed studies reveal a fundamental asymmetry in task difficulty. Structural tasks, such as tooth segmentation and enumeration, are relatively low-variance problems and achieve consistently high performance across architectures~\cite{bonfanti2025deep}. In contrast, disease detection—particularly early-stage caries—remains significantly more challenging due to subtle and high-variance visual features.

This disparity indicates that anatomical recognition and pathological detection impose fundamentally different requirements on model design and data quality. Even high-performing systems demonstrate reduced accuracy when transitioning from structural tasks to disease detection within the same framework~\cite{nassiri2025yolo}.

\subsubsection{Multi-Task Learning vs. Single-Task Models}

As shown in Table~\ref{tab:comparison}, the majority of studies adopt single-task formulations, with only a limited number exploring multi-task learning approaches. This represents a critical limitation, as real-world dental diagnosis requires simultaneous evaluation of multiple conditions.

Although a few studies attempt multi-task learning by combining diagnostic objectives within a single framework~\cite{chen2023deep}, such approaches remain rare and are often limited in scope or modality. The lack of integrated multi-task systems highlights a significant gap between current research and clinically deployable solutions.

\subsection{Deployment and Evaluation Challenges}

\subsubsection{Explainable AI}

Explainability is incorporated in only a small fraction of the reviewed studies (approximately 6.4\%), with most systems remaining black-box models~\cite{chen2023deep}. This lack of interpretability reduces clinical trust and poses challenges for regulatory approval, where transparency and human oversight are required. Consequently, explainable AI is not merely an enhancement but a necessary component for clinical deployment.

\subsubsection{Lack of Standardized Evaluation Metrics}

A major limitation across the reviewed literature is the lack of standardized evaluation metrics. Studies report heterogeneous measures, including mAP, Dice, IoU, accuracy, and F1-score, making direct comparison difficult. As reflected in Table~\ref{tab:comparison}, this inconsistency complicates benchmarking across models and tasks.

Systematic reviews confirm that metric variability contributes significantly to inter-study heterogeneity and limits reliable evidence synthesis~\cite{bonfanti2025deep}. The development of standardized evaluation protocols is therefore essential for advancing the field.

\subsection{Final Observation}

Overall, CNN-based and YOLO-driven approaches dominate current dental AI systems, while transformer-based models are emerging as a promising but resource-intensive alternative. Structural tasks achieve high reliability, whereas disease detection—particularly early-stage and quantitative assessment—remains significantly more challenging. These trends highlight a critical disconnect between current research, which focuses on isolated high-performance tasks, and the requirements of clinically deployable systems that demand integrated, explainable, and generalizable solutions across the full diagnostic pipeline.

% -----------------------------------------------------------------------
\section{Critical Analysis of Existing Methods}
\label{sec:critical}

\subsection{Metrics and Evaluation Limitations}

A major limitation across the reviewed literature is the lack of
standardized evaluation practices. Studies report a wide range of
metrics -- including mAP, Dice coefficient, accuracy, AUC, and
F1-score -- often with inconsistent thresholds and experimental setups.
This heterogeneity makes direct comparison across methods difficult and
obscures the identification of genuinely superior approaches
\cite{iacob2025automated,bonfanti2025deep,khubrani2025detection}.

Furthermore, many studies report high performance under controlled
experimental conditions without standardized benchmarking protocols.
Only a limited subset of works adopt consistent evaluation settings or
cross-study comparable methodologies.

Consequently, the absence of unified benchmarks and evaluation standards
fundamentally limits the reproducibility, comparability, and reliability
of existing approaches.

\subsection{Dataset Bias and Generalization Issues}

A critical issue in current research is the heavy reliance on
single-center datasets. A significant proportion of the reviewed studies
utilize data collected from a single institution or imaging setup, which
fails to capture variability across populations, imaging devices, and
acquisition conditions \cite{albano2024artificial,luke2025accuracy}.

This limitation leads to substantial generalization gaps. Models trained
on such datasets often demonstrate strong internal performance but
degrade significantly when evaluated on external data, as reflected in
multiple studies reporting notable performance drops under cross-dataset
validation \cite{muresanu2024automating}.

In addition, publicly available datasets remain limited in both scale
and diversity. Existing benchmarks such as DENTEX and CariesXrays
provide useful baselines but do not fully represent multi-disease,
multi-center clinical variability \cite{nassiri2025yolo,chen2024cariesxrays}.

Overall, the lack of diverse, standardized, multi-institutional datasets
and consistent validation protocols fundamentally constrains the
generalizability of current systems, thereby limiting their reliability
in real-world clinical deployment.

\subsection{Clinical Integration Gap}

Despite strong performance in controlled settings, most systems remain
confined to offline evaluation and focus on producing intermediate
outputs -- such as bounding boxes, segmentation masks, or classification
labels -- without integration into real-time clinical decision support.

Moreover, clinician involvement in model development and evaluation is
limited, and only a small number of studies assess real-world impact in
terms of diagnostic accuracy, treatment planning, or workflow efficiency
\cite{schwendicke2021cost,lee2025evaluation}.

Additionally, the limited incorporation of explainable AI further
restricts clinical trust, as most systems operate as black-box models
without interpretable outputs. This results in a clear disconnect
between technical performance and clinical applicability, indicating
that current systems are not yet aligned with real-world diagnostic
requirements.

\subsection{Fragmentation of Existing Approaches}

Another fundamental limitation is the fragmented nature of existing
methods. As evidenced in Table~\ref{tab:taxonomy}, the majority of
reviewed studies focus on isolated tasks such as detection, segmentation,
or classification, with only a small fraction exploring integrated
multi-task frameworks.

However, real-world dental diagnosis requires simultaneous evaluation of
multiple conditions, including caries detection, periodontal assessment,
severity staging, and quantitative measurement. The absence of systems
that jointly address these components prevents current approaches from
reflecting the complexity of clinical workflows.

To the best of our knowledge, no reviewed system operating on panoramic
radiographs simultaneously integrates detection, localization, severity
staging, quantification, and explainability within a unified framework.
This fragmentation represents a structural limitation of the field
rather than a limitation of individual model performance alone.

Addressing this limitation requires a shift from task-specific
optimization toward system-level design, where multiple diagnostic
components are jointly modeled and aligned with clinical requirements.

\vspace{4pt}
\noindent\textit{These limitations collectively highlight the need for
unified, clinically aligned, and generalizable dental AI systems, which
are discussed in the following section.}

% -----------------------------------------------------------------------
\section{Cross-Cutting Gap Analysis}
\label{sec:gaps}

Synthesizing findings across all 47 reviewed papers, five critical
interconnected gaps emerge that consistently limit clinical translational
potential.

\subsection{Gap 1: Single-Disease Focus}

The dominant pattern across all categories is single-disease training.
Of 47~papers reviewed, only Chen et al.~\cite{chen2023deep}
simultaneously detect both caries and periodontal disease -- exclusively
on periapical, not panoramic, radiographs. On panoramic radiographs
specifically, no reviewed study provides simultaneous caries detection,
severity staging, and periodontal bone loss quantification within a
single pipeline. This is the primary architectural gap relative to
clinical reality, where patients routinely present with co-occurring
conditions requiring simultaneous assessment from one exposure.

\subsection{Gap 2: Absence of Automated Severity Staging}

Of 18 C1~papers, only Asadi~\cite{asadi2024detection} and Bizjak and
Robic~\cite{bizjak2024advancing} address any severity distinction beyond
binary caries presence. In C2, only Wimalasiri et
al.~\cite{wimalasiri2025ai} address bone loss severity patterns. The
2018 AAP/EFP classification~\cite{chmielewski20252018} defines four
periodontitis stages, none addressed by any panoramic OPG system
reviewed. Severity staging is not optional -- it determines treatment
planning, referral, and prognosis. A system that detects caries without
staging it provides limited clinical decision support value.

\subsection{Gap 3: Explainability at 6.4\%}

Only 3 of 47~papers implement any XAI output. The consensus that
explainability is necessary for clinical adoption
\cite{nassiri2025yolo,gamal2025automated,sastika2025ai,ferrara2025applications}
is not reflected in implementation rates. Regulatory frameworks including
the EU AI Act classify AI diagnostic tools as high-risk systems requiring
transparency and human oversight -- requirements that current black-box
dental AI systems cannot satisfy. For junior clinicians specifically,
explainability is not optional; it is the mechanism through which
AI-assisted diagnosis becomes educational rather than merely predictive.

\subsection{Gap 4: Dataset Fragmentation}

Approximately 60\% of reviewed primary studies use private,
single-center datasets. Systematic reviews
\cite{iacob2025automated,albano2024artificial,luke2025accuracy,bonfanti2025deep,khubrani2025detection}
consistently report $I^2\,=\,88$--$96\%$ heterogeneity as the primary
barrier to evidence synthesis. The quantified real cost is documented by
Mure\c{s}anu et al.~\cite{muresanu2024automating}: F1 drops 22\% from
internal to external validation. The field requires a single
large-scale, multi-disease, multi-center, publicly available panoramic
OPG dataset with standardized caries and periodontal annotations --
nothing in the reviewed literature currently provides this.

\subsection{Gap 5: Clinical Workflow Integration}

With the exception of Schwendicke et al.~\cite{schwendicke2021cost} and
Lee et al.~\cite{lee2025evaluation}, no reviewed study evaluates AI
within an actual clinical workflow or measures clinician acceptance and
diagnostic impact. The majority of systems produce bounding boxes or
segmentation masks without clinical report generation, severity
summaries, or treatment recommendations. The economic evidence from
Schwendicke et al.~\cite{schwendicke2021cost} demonstrating
ICER\,=\,$-$13.9~\euro/year for AI-assisted caries detection represents
an isolated example of the real-world impact analysis the field urgently
needs.

% -----------------------------------------------------------------------
\section{Research Roadmap Toward Unified and Clinically Deployable
Dental AI Systems}
\label{sec:roadmap}

Building on the limitations identified in Section~\ref{sec:critical},
this section outlines a structured roadmap toward the development of
unified, clinically deployable dental AI systems.

\subsection{Short-Term Goals (1--2 Years)}

\textbf{Public Multi-Disease Panoramic Dataset:} The most immediate
bottleneck is the absence of a large-scale, standardized panoramic
dataset with multi-disease annotations. A consortium-driven effort is
required to develop a dataset containing at least 5,000 panoramic
radiographs with consistent annotations for caries (per-tooth with
severity), periodontal bone loss (CEJ--bone crest distance), and tooth
numbering (FDI standard).

Existing datasets such as CariesXrays \cite{chen2024cariesxrays} and
DENTEX \cite{nassiri2025yolo} provide partial solutions; however, they
lack comprehensive multi-disease annotations. Extending and integrating
these datasets represents a practical short-term direction.

\textbf{Explainability as a Standard Component:} Explainability must be
incorporated as a standard component using techniques such as Grad-CAM
or attention visualization, with validation against expert annotations
to ensure anatomical relevance
\cite{kaushik2025densenet201,bilal2024pdcnet}.

\subsection{Mid-Term Goals (3--5 Years)}

\textbf{Unified Multi-Disease System for Panoramic Radiographs:} The
most critical research direction is the development of a unified
multi-disease system capable of simultaneously detecting caries,
performing severity staging, quantifying periodontal bone loss, and
supporting comprehensive multi-condition diagnosis from a single
panoramic radiograph.

This requires multi-head architectures with shared feature
representations and task-specific outputs. While individual components
exist in the literature, their integration into a single coherent
framework remains an open challenge
\cite{nassiri2025yolo,chen2023deep,jundaeng2025advanced,
muresanu2024automating}.

\textbf{Clinical Alignment with Standardized Diagnostic Frameworks:}
Future systems must generate outputs aligned with established clinical
classification standards, such as ICCMS for caries and the 2018 AAP/EFP
classification for periodontal disease \cite{chmielewski20252018}. Such
alignment transforms model outputs into clinically interpretable
results, thereby bridging the gap between algorithmic predictions and
clinically actionable decisions.

\subsection{Long-Term Goals (Clinical Deployment)}

\textbf{Prospective Clinical Validation and Deployment:} Future
research must move beyond retrospective evaluation toward prospective
clinical validation. This includes randomized studies assessing
diagnostic accuracy, clinician efficiency, and patient outcomes with
and without AI assistance. Such studies should follow standardized
reporting protocols to ensure reproducibility and comparability across
clinical settings.

The cost-effectiveness framework proposed by Schwendicke et al.\
\cite{schwendicke2021cost} provides a foundation for evaluating
clinical impact and should be extended to multi-disease systems and
panoramic radiographs. Such validation is essential not only for
regulatory approval but also for establishing trust among clinicians
and enabling widespread adoption.

\vspace{4pt}
\noindent\textit{\textbf{Final Perspective:} Collectively, this roadmap
outlines a progression from data standardization and model
explainability toward fully integrated, clinically validated systems,
providing a clear and actionable pathway for the next generation of
dental AI research.}
% -----------------------------------------------------------------------
\section{Conclusion}
\label{sec:conclusion}

This survey reviewed 47 deep learning studies (2021--2026) on dental
caries detection, periodontal bone loss assessment, and related tasks
from panoramic and bitewing dental radiographs. The literature
demonstrates that deep learning has achieved clinically competitive
performance on individual sub-tasks: caries detection reaches 96\%
recall with Attention U-Net \cite{pornprasertsuk2025clinical};
periodontal landmark detection achieves mAP50\,=\,0.995 with YOLOv8
\cite{jundaeng2025advanced}; dual-disease classification attains
AUC\,$>$\,98\% on periapical X-rays with EfficientNet-B0 and Grad-CAM
\cite{chen2023deep}; and AI-assisted caries detection is economically
superior to unaided clinical assessment with
ICER\,=\,$-$13.9~\euro/year \cite{schwendicke2021cost}.

However, five structural gaps persist across the entire reviewed corpus:
single-disease models that ignore clinical co-occurrence; near-absence of
automated severity staging despite its treatment planning centrality; a
6.4\% explainability implementation rate despite universal consensus on
its necessity; dataset fragmentation from private single-center
collections; and minimal clinical workflow integration beyond isolated
studies.

No existing system simultaneously addresses caries detection, severity
staging, periodontal bone loss quantification, and explainability on
panoramic radiographs within a single unified pipeline. This is the
definitive unmet need that the reviewed literature collectively
establishes. The technical components exist -- YOLO detection, U-Net
segmentation, Grad-CAM explainability, EfficientNet classification, and
transformer-based global attention have all been individually
demonstrated on dental radiographic data. The challenge is integration,
clinical validation, and standardization toward systems that are not
merely accurate but clinically trustworthy, interpretable, and
actionable -- particularly for the junior clinicians and dental students
who stand to benefit most from AI-assisted diagnostic support.

% -----------------------------------------------------------------------
\section*{Acknowledgment}
The authors thank Dr.\ Mamatha H R, Director, Bajaj Engineering Skills
Training@PESU, PES University -- RR Campus, Bangalore, for her guidance
and supervision throughout this work.

% -----------------------------------------------------------------------
\bibliographystyle{IEEEtran}
\bibliography{source}

\end{document}